\documentclass[12pt]{article}
\usepackage[margin=0.72in]{geometry}
\usepackage{lmodern}
\usepackage{microtype}
\usepackage{multicol}
\usepackage{fancyhdr}
\usepackage{titlesec}
\usepackage{enumitem}
\usepackage{xcolor}
\usepackage[hidelinks]{hyperref}

\setlength{\columnsep}{0.30in}
\setlength{\parindent}{1.05em}
\setlength{\parskip}{0.24em}
\setlength{\headheight}{15pt}
\titleformat{\section}{\large\bfseries\scshape}{}{0pt}{}
\titleformat{\subsection}{\normalsize\bfseries}{}{0pt}{}
\pagestyle{fancy}
\fancyhf{}
\fancyhead[L]{Whately: Making Clear Claims}
\fancyhead[R]{Reading Handout}
\fancyfoot[C]{\thepage}
\renewcommand{\headrulewidth}{0.4pt}

\newcommand{\focusbox}[1]{%
\vspace{0.4em}\noindent\colorbox{gray!12}{\begin{minipage}{0.96\linewidth}#1\end{minipage}}\vspace{0.5em}}

\begin{document}
\begin{center}
{\LARGE\bfseries What We Say and What We Mean}\par\vspace{0.15em}
{\large\scshape Reading 3: Propositions and Clear Claims}\par\vspace{0.45em}
\rule{0.82\textwidth}{0.4pt}\par\vspace{0.6em}
\end{center}

\section*{Vocabulary Preview}
\begin{multicols}{2}
\begin{itemize}[leftmargin=*, itemsep=0.15em]
\item \textbf{Proposition}: A statement that can be true or false.
\item \textbf{Subject}: What the statement is about.
\item \textbf{Predicate}: What is said about the subject.
\item \textbf{Quality}: Whether a statement affirms or denies something (affirmative or negative).
\item \textbf{Quantity}: How much of the subject is included (all, some, or none).
\item \textbf{Universal}: A statement that applies to every member of a group (All S are P, No S are P).
\item \textbf{Particular}: A statement that applies to only some members of a group (Some S are P, Some S are not P).
\item \textbf{Distribution}: Whether a term refers to all or only part of the group it names.
\item \textbf{Copula}: The linking word (usually “is,” “are,” “is not,” or “are not”).
\end{itemize}
\end{multicols}

\section*{Introduction}
In Unit 2 we learned that words can shift their meaning and break arguments. In this unit we move one step further. Even when the meaning of our terms is fixed, we can still make unclear or misleading claims if we do not pay attention to how we combine terms into statements.

Richard Whately emphasized that logic is concerned with the structure of reasoning. Before we can evaluate whether one statement follows from another, we must first be able to see exactly what a statement is claiming. This requires understanding the difference between the subject and predicate, whether a claim is affirmative or negative, whether it applies to all or only some members of a group, and whether the terms are being used in a distributed or undistributed way.

This reading introduces the basic structure of propositions and the most common ways claims can be stated clearly or unclearly.

\focusbox{\textbf{Source note:} This reading draws on the treatment of propositions found in Book II of Whately’s \textit{Elements of Logic}, with additional examples and explanations added for clarity.}

\begin{multicols}{2}
\section*{Reading}

\subsection*{1. What Is a Proposition?}
A \textit{proposition} is a statement that can be true or false. It is not a question, a command, or an exclamation. It is a claim about the world that someone could agree with or disagree with.

Examples of propositions:
\begin{itemize}
\item All humans are mortal.
\item Some politicians are honest.
\item No triangle has four sides.
\item Socrates is wise.
\end{itemize}

These statements can each be evaluated as true or false. By contrast, “Is Socrates wise?” and “Be wise!” are not propositions.

\subsection*{2. Subject and Predicate}
Every proposition has two main parts. The \textit{subject} is what the statement is about. The \textit{predicate} is what is said about the subject. These two parts are joined by a linking word called the \textit{copula} (usually a form of “to be”).

Consider the statement: “All students are responsible.”

- Subject: students
- Predicate: responsible
- Copula: are

The subject tells us what we are talking about. The predicate tells us what we are saying about it.

\subsection*{3. Quality: Affirmative and Negative}
Propositions differ in \textit{quality}. A proposition is \textit{affirmative} when it asserts that the predicate belongs to the subject. It is \textit{negative} when it denies that the predicate belongs to the subject.

Affirmative examples:
\begin{itemize}
\item All birds are animals.
\item Some metals are magnetic.
\end{itemize}

Negative examples:
\begin{itemize}
\item No birds are mammals.
\item Some metals are not magnetic.
\end{itemize}

The difference matters because an affirmative claim and a negative claim have very different logical consequences. Saying “All S are P” is not the same as saying “No S are P.”

\subsection*{4. Quantity: Universal and Particular}
Propositions also differ in \textit{quantity}. A proposition is \textit{universal} when it makes a claim about every member of the subject class. It is \textit{particular} when it makes a claim about only some members of the subject class.

Universal propositions:
\begin{itemize}
\item All humans are mortal. (Every human is mortal.)
\item No humans are immortal. (Not a single human is immortal.)
\end{itemize}

Particular propositions:
\begin{itemize}
\item Some humans are philosophers. (At least one human is a philosopher.)
\item Some humans are not philosophers. (At least one human is not a philosopher.)
\end{itemize}

Note that the word “some” in logic means “at least one” and does not exclude the possibility that all members satisfy the predicate. This is different from everyday usage, where “some” often implies “not all.”

\subsection*{5. The Four Standard Forms}
When we combine quality and quantity, we get four standard forms of propositions. These are traditionally labeled with the letters A, E, I, and O.

\begin{itemize}
\item \textbf{A}: All S are P (universal affirmative)
\item \textbf{E}: No S are P (universal negative)
\item \textbf{I}: Some S are P (particular affirmative)
\item \textbf{O}: Some S are not P (particular negative)
\end{itemize}

These four forms are the basic building blocks of deductive reasoning. Most of the claims we make in ordinary language can be translated into one of these four forms once we identify the subject, predicate, quality, and quantity.

\subsection*{6. Distribution of Terms}
A term is said to be \textit{distributed} when the proposition makes a claim about every member of the class the term names. It is \textit{undistributed} when the proposition makes a claim about only some members of the class.

Consider the four forms again:

- In an A proposition (“All S are P”), the subject term S is distributed, but the predicate term P is not. We are saying something about every S, but we are not saying that every P is an S.
- In an E proposition (“No S are P”), both S and P are distributed. We are saying that no member of S is a member of P, and no member of P is a member of S.
- In an I proposition (“Some S are P”), neither term is distributed. We are only making a claim about some members of each class.
- In an O proposition (“Some S are not P”), only the subject S is distributed in a limited way; the predicate is not.

Understanding distribution is essential for determining whether a conclusion follows from premises in a syllogism. This will become especially important in later units.

\subsection*{7. Why These Distinctions Matter}
Many arguments become unclear or invalid because the speaker has not made the quantity or quality of their claims explicit. Consider the following exchange:

Person A: “Politicians are dishonest.”
Person B: “That is not true. My representative is very honest.”

The problem here is that Person A has not made clear whether they meant “All politicians are dishonest” or “Some politicians are dishonest.” If they meant the first, Person B has offered a counterexample. If they meant the second, Person B has not contradicted the claim at all.

Clear reasoning requires that we know exactly what is being asserted about how much of the subject and in what way.

\subsection*{8. Translating Ordinary Language}
In everyday speech, people rarely use the precise forms “All S are P” or “Some S are not P.” Instead, they use a wide variety of expressions. Part of clear thinking involves translating these varied expressions into their logical form.

Examples of translation:
\begin{itemize}
\item “Every student passed the test” → All students are people who passed the test.
\item “No one likes that policy” → No people are people who like that policy.
\item “Most birds can fly” → Some birds are animals that can fly. (Note: “most” is usually treated as “some” in basic logic.)
\item “Only citizens may vote” → All people who may vote are citizens.
\item “None but the brave deserve the fair” → All people who deserve the fair are brave people.
\end{itemize}

Learning to recognize these patterns helps students move from the vagueness of ordinary language to the precision required for careful argument.

\section*{Reading Focus}
\begin{enumerate}[leftmargin=*]
\item What is the difference between the subject and the predicate of a proposition?
\item Explain the difference between quality and quantity in a proposition. Give one example of each.
\item Why is the word “some” in logic understood to mean “at least one” rather than “some but not all”?
\item In an A proposition (“All S are P”), which term is distributed and which is undistributed? Why does this distinction matter?
\item Why can the statement “Politicians are dishonest” lead to misunderstanding in an argument?
\end{enumerate}

\section*{Window Note and Summary}
Create a window note that shows one proposition written in four different ways by changing only its quality and quantity (A, E, I, and O forms). For each version, include a small drawing or symbol that represents what the claim is asserting. At the bottom of the window note, write a 5–7 sentence summary explaining why being precise about quality and quantity helps us evaluate whether one claim follows from another.

\section*{Connection to Unit 2}
How does the ability to identify the exact quantity and quality of a claim help us avoid the kind of ambiguity we studied in Unit 2 (when a word shifts meaning)?

\section*{Bridge to the Next Reading}
In the next unit we will begin combining propositions into arguments called syllogisms. Before we can judge whether a conclusion follows from two premises, we must first be able to see clearly what each premise is actually claiming about its subject. The skills practiced in this unit will be essential for that work.

\end{multicols}
\end{document}